\section{В чем задача и почему ломается бейзлайн}

Когда я только начал разбираться с Forward-Backward (FB) представлениями, идея показалась мне потрясающе элегантной. Вместо того чтобы под каждую новую задачу заново учить RL-агента с нуля, мы один раз скармливаем оффлайн-датасет модели и раскладываем меру будущих посещений (Successor Measure) на две половинки:
\addequation{
    M^{\pi_z}(s, s') = \mathbb{E}_{\pi_z} \left[ \sum_{t=0}^\infty \gamma^t \mathbb{I}(s_t = s') \;\middle|\; s_0 = s \right] \approx F(s, z)^\top B(s').
}{eq:fb_factorization}

Физический смысл скалярного произведения $F(s, z)^\top B(g)$ очень наглядный: это не геометрическое расстояние, а оценка того, насколько вероятно роботу добраться до цели $g$ из текущей точки $s$, если он выберет латентную интенцию $z \in \mathbb{S}^{d-1}$.

\begin{figure}[H]
    \centering
    \begin{tikzpicture}[>=stealth, font=\small, scale=0.95, every node/.style={transform shape}]
    % Левая часть: Single-Intention Baseline
    \begin{scope}[shift={(0,0)}]
        \draw[thick, fill=gray!10] (0,0) rectangle (5,3.2);
        \draw[ultra thick, line width=3pt, red!70!black] (2.2, 0.6) -- (2.2, 2.6);
        \node[font=\footnotesize\bfseries, red!70!black] at (2.2, 2.85) {Стена};
        
        \node[circle, fill=green!70!black, inner sep=2.5pt, label=below:{\footnotesize Старт $s_0$}] (S) at (0.8, 1.6) {};
        \node[star, star points=5, star point ratio=2.25, fill=yellow!80!black, draw=black, inner sep=2pt, label=below:{\footnotesize Цель $g$}] (G) at (4.2, 1.6) {};
        
        \draw[->, ultra thick, dashed, red!80!black] (S) -- (2.05, 1.6);
        \node[cross out, draw=red, line width=2pt, minimum size=10pt] at (2.15, 1.6) {};
        \node[fill=white, inner sep=2pt, font=\scriptsize, align=center] at (1.1, 2.4) {$\pi^h(s_0, B(g)) \to z_{\text{direct}}$\\(Застревание в стене)};
        
        \node[font=\bfseries, align=center] at (2.5, -0.4) {а) Single-Intention Baseline};
    \end{scope}

    % Правая часть: Multi-Subgoal Sequence
    \begin{scope}[shift={(6.5,0)}]
        \draw[thick, fill=gray!10] (0,0) rectangle (5,3.2);
        \draw[ultra thick, line width=3pt, red!70!black] (2.2, 0.6) -- (2.2, 2.6);
        \node[font=\footnotesize\bfseries, red!70!black] at (2.2, 2.85) {Стена};
        
        \node[circle, fill=green!70!black, inner sep=2.5pt, label=below:{\footnotesize Старт $s_0$}] (S2) at (0.8, 1.6) {};
        \node[circle, fill=blue!70!black, inner sep=2pt, label=above:{\footnotesize $w_1$}] (W1) at (1.2, 2.7) {};
        \node[circle, fill=blue!70!black, inner sep=2pt, label=above:{\footnotesize $w_2$}] (W2) at (3.2, 2.7) {};
        \node[star, star points=5, star point ratio=2.25, fill=yellow!80!black, draw=black, inner sep=2pt, label=below:{\footnotesize Цель $g$}] (G2) at (4.2, 1.6) {};
        
        \draw[->, thick, blue!80!black] (S2) to[bend left=15] (W1);
        \draw[->, thick, blue!80!black] (W1) to (W2);
        \draw[->, thick, blue!80!black] (W2) to[bend left=15] (G2);
        
        \node[fill=white, inner sep=2pt, font=\scriptsize, align=center] at (3.8, 2.3) {Последовательность\\вейпоинтов};
        
        \node[font=\bfseries, align=center] at (2.5, -0.4) {б) Multi-Subgoal Planning};
    \end{scope}
\end{tikzpicture}

    \caption{Почему ломается бейзлайн: единая прямая интенция тянет робота прямо в стену (слева), тогда как цепочка промежуточных подцелей позволяет обойти тупик (справа).}
    \label{fig:single_vs_multi}
\end{figure}

В базовой архитектуре всё устроено в два уровня:
\begin{itemize}[leftmargin=*,noitemsep]
    \item \textbf{Low-level исполнитель} $\pi^\ell(a \mid s, z)$ — 8-суставный контроллер муравья в MuJoCo, отрабатывающий команду $z$.
    \item \textbf{High-level контроллер} $\pi^h(s, B(g)) \to z$ (в коде — \texttt{high\_actor}), выдающий одну глобальную команду на всю цель целиком.
\end{itemize}

И вот тут кроется главная засада. В открытом пространстве всё работает здорово. Но в лабиринте AntMaze муравей — это не невесомая 2D-точка, а тяжелый робот с физикой контактов и моментами инерции. Когда цель находится за углом или в соседнем коридоре, \texttt{high\_actor} пытается вести робота напрямую через стену (эффект «мухи на стекле»). Муравей втыкается в угол, теряет баланс, переворачивается и застревает. 

В итоге бейзлайн на лабиринте Medium выдает около 76\%, а на сложном Large полностью проваливается до 24\%. Моей целью стало научить агента \textbf{выстраивать цепочки промежуточных интенций (Multi-Subgoal Planning)}, не прибегая к тяжелым генеративным диффузионным моделям. Ниже я подробно расскажу, как развивались мои идеи.
